\documentclass[11pt,a4paper]{article}
\usepackage[T1]{fontenc}
\usepackage[utf8]{inputenc}
\usepackage[margin=25mm,headheight=15pt]{geometry}
\usepackage{amsmath,amsthm,mathtools}
\usepackage{newtxtext,newtxmath}
\usepackage{microtype,booktabs,longtable,array,enumitem}
\usepackage{xcolor,listings,fancyhdr,hyperref,xurl}
\usepackage[nameinlink,noabbrev]{cleveref}
\hypersetup{colorlinks=true,linkcolor=blue!45!black,citecolor=blue!45!black,urlcolor=blue!45!black,
 pdftitle={Real Inverse Asymptotics at a Singular Endpoint},
 pdfauthor={OpenAI, prepared for Vladimir Reshetnikov}}
\definecolor{codebg}{gray}{0.965}
\lstdefinestyle{wl}{basicstyle=\ttfamily\small,breaklines=true,columns=fullflexible,
 keepspaces=true,showstringspaces=false,backgroundcolor=\color{codebg},
 frame=single,rulecolor=\color{black!15},framesep=5pt,
 morecomment=[s]{(*}{*)},commentstyle=\color{black!60},
 literate={->}{{$\to$}}2}
\lstset{style=wl}
\pagestyle{fancy}\fancyhf{}
\fancyhead[L]{\small Real inverse asymptotics}
\fancyhead[R]{\small Theory and implementation}
\fancyfoot[C]{\thepage}
\setlength{\parskip}{4pt}
\setlength{\parindent}{1.25em}
\setlist{nosep,leftmargin=1.6em}
\numberwithin{equation}{section}
\newtheorem{theorem}{Theorem}[section]
\newtheorem{lemma}[theorem]{Lemma}
\newtheorem{proposition}[theorem]{Proposition}
\newtheorem{corollary}[theorem]{Corollary}
\theoremstyle{definition}
\newtheorem{definition}[theorem]{Definition}
\newtheorem{example}[theorem]{Example}
\newtheorem{remark}[theorem]{Remark}
\newcommand{\R}{\mathbb R}
\newcommand{\C}{\mathbb C}
\newcommand{\N}{\mathbb N}
\newcommand{\OO}{\mathcal O}
\newcommand{\dd}{\mathrm d}
\newcommand{\DL}{\partial_L}
\newcommand{\nn}{\boldsymbol n}
\newcommand{\ee}{\boldsymbol e}
\newcommand{\del}{\boldsymbol\delta}
\newcommand{\PP}{\mathcal P}
\newcommand{\IA}{\mathcal I_A}
\newcommand{\BA}{\mathcal B_A}
\newcommand{\fall}[2]{\left(#1\right)^{\underline{#2}}}
\newcommand{\rise}[2]{\left(#1\right)^{\overline{#2}}}
\newcommand{\code}[1]{\texttt{#1}}
\DeclareMathOperator{\val}{val}
\DeclareMathOperator{\Res}{Res}
\DeclareMathOperator{\supp}{supp}

\begin{document}
\hypersetup{pageanchor=false}
\begin{titlepage}
\vspace*{12mm}
\noindent\textsc{Constructive asymptotic analysis \quad / \quad Wolfram Language}
\vspace{18mm}
\begin{flushleft}
{\Huge\bfseries Real Inverse Asymptotics\\[4pt]at a Singular Endpoint}\par
\vspace{8mm}
{\Large Power--logarithmic expansions, irrational exponents,\\
branch selection, and a finite inversion algorithm}\par
\vspace{15mm}
{\large A research article and implementation prepared for\\
Vladimir Reshetnikov}\par
\vspace{5mm}
{\large September 7, 2026 \quad\textbullet\quad Package version 1.0.0}
\end{flushleft}
\vfill
\noindent\begin{minipage}{0.94\textwidth}
\textbf{The principal result.}
For $f(x)=c x^p(1+\sum_{j=1}^r x^{\delta_j}P_j(\log x))$, with
$c,p,\delta_j>0$, the positive inverse has explicitly computable polynomial
coefficients on the locally finite support $1+\N\del$. The construction
works without analyticity at zero, merges coincident exponents exactly,
and supplies a logarithm-aware remainder at every power cutoff.

\medskip
\textbf{Validation boundary.}
Seventy-nine independent symbolic and 250-digit numerical checks were run
successfully in Python. The Wolfram package and a 33-test native suite are
included. No native Wolfram execution is claimed: the connected evaluator
returned an unavailable-endpoint error in this session. Formal residual
checks in the package are not numerical interval certificates.
\end{minipage}
\end{titlepage}
\hypersetup{pageanchor=true}
\pagenumbering{roman}

\begin{abstract}
The inverse of $x+x^2(1+\log x)$ at $0^+$ is not an ordinary Taylor series,
and the inverse of $x+x^{\sqrt2}$ requires irrational exponents. Both are
instances of an elementary but useful inversion calculus for finite
power--logarithmic perturbations of a positive monomial. We prove existence
and selection of the positive real branch, a closed multi-index coefficient
formula, local convergence of the resulting generalized series for exact
finite models, and quantitative asymptotic remainder orders for arbitrary
weighted truncations. A finite border of an order ideal determines the first
possible omitted power and a safe logarithmic degree, even when generator
weights collide. A separate Newton algorithm in a truncated logarithmic
jet algebra gives the same inverse, and forward residuals transfer to
inverse errors after the correct derivative rescaling. We also treat
finite forward expansions with differentiated remainder bounds, explain
why a bound on the forward function alone need not preserve invertibility,
and state the boundaries of the implemented class. The accompanying
Wolfram Language package contains both coefficient and Newton engines,
exact exponent ordering, explicit failure cases, residual inspection, and a
more general auxiliary-parameter inversion utility. Complete examples,
proofs, reproducible verification programs, a native regression suite,
and the full package listing accompany the article.
\end{abstract}

\tableofcontents
\newpage
\pagenumbering{arabic}

\section{The problem and the answer}
\label{sec:answer}

\subsection{The two examples}
Vladimir Reshetnikov's question, posted on December 11, 2020 and later
updated with an irrational-power example, asks for the real inverse near
zero of
\begin{equation}\label{eq:original}
 f(x)=x+x^2(1+\log x),\qquad x>0,
\end{equation}
and for an analogous expansion of the inverse of
\begin{equation}\label{eq:irr-original}
 f_\alpha(x)=x+x^\alpha,\qquad \alpha=\sqrt2.
\end{equation}
The uploaded Markdown copy and the public question agree on these two
problems \cite{question}. The reported computer-algebra expressions in that
question are historical observations, not tests of the present version of
Mathematica.

Write $L=\log y$. The first answer begins
\begin{align}
 g(y)={}&y-y^2(1+L)+y^3(2L^2+5L+3)\notag\\
 &-y^4\left(5L^3+\frac{41}{2}L^2+27L+\frac{23}{2}\right)\notag\\
 &+y^5\left(14L^4+\frac{241}{3}L^3+\frac{335}{2}L^2
                  +151L+\frac{299}{6}\right)\notag\\
 &+\OO\!\left(y^6(1+|L|)^5\right),\qquad y\longrightarrow0^+.
 \label{eq:log-answer}
\end{align}
In particular, the correct two-block statement is
\begin{equation}\label{eq:two-blocks}
 g(y)=y-y^2(1+\log y)+\OO\!\left(y^3(1+|\log y|)^2\right).
\end{equation}
It is \emph{not} valid to replace this remainder by the ordinary analytic
Landau bound $\OO(y^3)$. Indeed, the next coefficient is
$2L^2+5L+3$, which is unbounded. One may instead write
$\OO(y^{3-\varepsilon})$ for every fixed $\varepsilon>0$, but that loses
useful information.

For the second example,
\begin{align}
 g_\alpha(y)
 &=\sum_{n=0}^{\infty} b_n(\alpha)y^{1+n(\alpha-1)},\label{eq:alpha-series}\\
 b_n(\alpha)
 &=\frac{(-1)^n}{n!}\fall{\alpha n}{n-1}\quad(n\ge1),
 \qquad b_0(\alpha)=1.\label{eq:alpha-coeff}
\end{align}
At $\alpha=\sqrt2$, this gives exactly
\begin{align}
 g_{\sqrt2}(y)={}&y-y^{\sqrt2}+\sqrt2\,y^{2\sqrt2-1}
 -\frac{6-\sqrt2}{2}y^{3\sqrt2-2}\notag\\
 &+\OO\!\left(y^{4\sqrt2-3}\right).
 \label{eq:irr-answer}
\end{align}
Both formulas follow from the same theorem below. In fact, for the exact
finite models considered here, sufficiently small positive $y$ gives
convergent generalized series, not merely formal ones.

\subsection{Immediate use of the package}
Extract the distribution and evaluate the following, with the symbols
\code{x}, \code{y}, and \code{z} unassigned:
\begin{lstlisting}
Get["path/to/real-inverse-series/Kernel/RealInverseSeries.wl"];
Clear[x, y, z];

r1 = RealInverseSeries[x + x^2 (1 + Log[x]), {x, y}, 6];
r1["Expansion"]
r1["Remainder"]

r2 = RealInverseSeries[x + x^Sqrt[2], {x, z}, 4 Sqrt[2] - 3];
r2["Expansion"]
r2["Remainder"]
\end{lstlisting}
The third argument is a \emph{strict power cutoff}: retain all complete
logarithmic blocks whose exponent in the output variable is smaller than
that number. It is not a count of monomials and not a count of
auxiliary-parameter orders. Thus cutoff $6$ in the first call retains the
five blocks displayed in \eqref{eq:log-answer}; cutoff $4\sqrt2-3$ in the
second call retains precisely the four terms in \eqref{eq:irr-answer}.
The returned association keeps the ordinary expression separate from its
mathematically valid remainder scale.

\section{Selecting the real branch before expanding}
\label{sec:branch}

\subsection{Global uniqueness in the original logarithmic example}
\begin{proposition}\label{prop:global-log}
The function in \eqref{eq:original} is a strictly increasing bijection
$(0,\infty)\to(0,\infty)$. Its inverse satisfies $g(y)\sim y$ at $0^+$.
It extends to a $C^1$ function at zero with $g(0)=0$ and $g'(0)=1$, but
has no $C^2$ extension there.
\end{proposition}
\begin{proof}
Direct differentiation gives
\[
 f'(x)=1+x(3+2\log x),\qquad f''(x)=5+2\log x.
\]
The function $x(3+2\log x)$ has its global minimum at
$x=e^{-5/2}$, with value $-2e^{-5/2}$. Consequently
\[
 f'(x)\ge 1-2e^{-5/2}>0\qquad(x>0).
\]
Also $f(0^+)=0$ and $f(x)\to\infty$ as $x\to\infty$. The intermediate
value theorem and strict monotonicity prove the bijection and positivity.
Since $f(x)/x=1+x(1+\log x)\to1$, inversion gives $g(y)/y\to1$.
For positive arguments, $g'(y)=1/f'(g(y))\to1$. Finally,
\[
 g''(y)=-\frac{f''(g(y))}{f'(g(y))^3}\longrightarrow+\infty.
\]
Thus the second derivative cannot extend continuously to zero.
\end{proof}

A complex inverse branch through a nonzero complex solution of $f(x)=0$
is a different local object. An expansion with a nonzero complex constant
term cannot represent the branch characterized by $g(y)>0$ and $g(y)\to0$.
The obstacle is not a missing choice of a Taylor coefficient: the correct
real branch is nonanalytic at the endpoint itself.

\subsection{The general finite model}
Our input class is
\begin{equation}\label{eq:model}
 f(x)=c x^p\left(1+R(x)\right),\qquad
 R(x)=\sum_{j=1}^r x^{\delta_j}P_j(\log x),
\end{equation}
where $c,p,\delta_j>0$ are fixed real numbers and each $P_j$ is a nonzero
real polynomial. Repeated values of $\delta_j$ may be merged at the input
stage. The case $r=0$ is the exact monomial inverse. Unless stated
otherwise, \eqref{eq:model} is an exact identity on a positive interval.

\begin{lemma}[Power beats logarithm]\label{lem:dominance}
For every $a>0$ and fixed real $b$, $x^a(1+|\log x|)^b\to0$ as
$x\to0^+$. In particular $R(x)=o(1)$ and $xR'(x)=o(1)$.
\end{lemma}
\begin{proof}
Put $x=e^{-t}$, so $t\to\infty$. The first assertion is
$e^{-at}(1+t)^b\to0$, which follows, for example, from repeated
l'H\^opital differentiation when $b$ is bounded above by an integer.
For a single term,
\[
 x\frac{\dd}{\dd x}\bigl(x^\delta P(\log x)\bigr)
 =x^\delta(\delta P+P')(\log x),
\]
and the same limit applies. There are only finitely many terms.
\end{proof}

\begin{theorem}[The positive inverse germ]\label{thm:branch}
For \eqref{eq:model}, there is $x_0>0$ such that $f$ is positive and
strictly increasing on $(0,x_0)$. Its inverse $g$ on $(0,f(x_0))$ is
unique there and satisfies
\begin{equation}\label{eq:scale}
 s=(y/c)^{1/p},\qquad L=\log s,\qquad \frac{g(y)}s\longrightarrow1.
\end{equation}
\end{theorem}
\begin{proof}
Positivity follows from $1+R(x)\to1$. Differentiation gives
\begin{equation}\label{eq:derivative-model}
 f'(x)=c p x^{p-1}\left[1+
 \sum_j x^{\delta_j}\left(
   \left(1+\frac{\delta_j}{p}\right)P_j(\log x)
       +\frac1pP_j'(\log x)\right)\right].
\end{equation}
The bracket tends to $1$, so the derivative is positive eventually. The
function tends to zero. The intermediate value theorem gives the inverse,
and its defining equation yields
$g(y)/s=(1+R(g(y)))^{-1/p}\to1$.
\end{proof}

The theorem selects a \emph{germ}: a function defined for sufficiently small
positive arguments, with two functions identified if they agree on a
smaller such interval. It does not assert global monotonicity of every
finite model. Negative correction coefficients are allowed; only the
leading coefficient and exponent must be positive.

\section{The algebra of complete logarithmic blocks}
\label{sec:algebra}

\subsection{Support and dominance}
Set
\begin{equation}\label{eq:semigroup}
 S=\{\nn\cdot\del:\nn\in\N^r\},\qquad
 \nn\cdot\del=\sum_j n_j\delta_j.
\end{equation}
Here $\N$ includes zero. A \emph{block} is $s^d Q(L)$ with
$Q\in\R[L]$. At zero, blocks of smaller power dominate blocks of larger
power, regardless of their fixed polynomial degrees. Within one block,
higher powers of $|L|$ dominate lower ones. It is usually clearer and less
error-prone to retain the whole polynomial at each exponent.

\begin{lemma}[Local finiteness]\label{lem:finite}
For every $A>0$, the set
\begin{equation}\label{eq:ideal}
 \IA=\{\nn\in\N^r:\nn\cdot\del<A\}
\end{equation}
is finite. Every element of $S$ has only finitely many multi-index
representations, and $S\cap[0,A)$ is finite.
\end{lemma}
\begin{proof}
Each $n_j<A/\delta_j$, so $\IA$ lies in a finite rectangular box.
For representations of a fixed weight $d$, use a box with $A>d$.
\end{proof}

It is important not to confuse this nonnegative semigroup with the group
$\sum_j\mathbb Z\delta_j$. An irrationally generated group may be dense,
but the positive semigroup here remains locally finite. Irrational
exponents therefore create no infinite computation below a finite cutoff.

If two different indices have the same weight, their polynomial
coefficients must be added. For example, gaps $1/2$ and $1$ satisfy
$2(1/2)=1$; the contributions indexed by $(2,0)$ and $(0,1)$ belong to the
same output block.

\subsection{A filtered algebra rather than a misleading Taylor series}
Consider formal sums
\begin{equation}\label{eq:algebra}
 U=\sum_{d\in S}s^d Q_d(L),\qquad Q_d\in\R[L].
\end{equation}
Addition is coefficientwise; multiplication uses the finite Cauchy sums
at every bounded weight. Terms of weight at least $A$ form an ideal, and
truncation modulo that ideal gives a finite logarithmic jet. A jet is
stored as a list of pairs $(d,Q_d)$, ordered by exact comparison of $d$.
No floating-point approximation of a radical exponent enters the ordering.
The valuation of a nonzero jet is its smallest weight with a nonzero
polynomial coefficient; the zero jet has valuation $+\infty$.

If $U$ has strictly positive support, then
\[
 (1+U)^a=\sum_{k\ge0}\binom ak U^k,
 \qquad
 \log(1+U)=\sum_{k\ge1}\frac{(-1)^{k+1}}k U^k
\]
are well-defined formal sums: below a fixed cutoff only finitely many
powers can contribute. A scalar $a\ne0$ plus a positive-order jet has a
reciprocal by the geometric series. We do not claim that all nonzero
polynomial-logarithmic leading blocks are units in this algebra.

The logarithmic coordinate is initially formal. Its connection to actual
functions is justified by the analytic theorem in \cref{sec:convergence},
not by pretending that $L$ is bounded as $s\to0$.

\begin{lemma}[Differential operator on a block]\label{lem:operator}
For real $a$ and polynomial $Q$,
\begin{equation}\label{eq:operator}
 \frac{\dd^k}{\dd s^k}\bigl(s^aQ(\log s)\bigr)
 =s^{a-k}\prod_{j=0}^{k-1}(a-j+\DL)Q(L).
\end{equation}
All factors commute.
\end{lemma}
\begin{proof}
The case $k=1$ is the product and chain rules. Applying the same identity
inductively reduces the power by one and appends the next operator.
Each factor is a polynomial in $\DL$, so the factors commute.
\end{proof}

The derivative term $\DL$ is essential. Treating $\log x$ as a frozen
coefficient when composing or inverting would omit the $5L+3$ correction
in the third block of \eqref{eq:log-answer}.

\section{Lagrange inversion at a moving positive point}
\label{sec:lagrange}

Classical Lagrange inversion has many equivalent forms; see
\cite{dlmf,gessel,surya}. For this problem the useful version is an
expansion in an \emph{auxiliary perturbation parameter}, at each fixed
positive $y$, not a Taylor expansion of the inverse at the singular point
$y=0$. This distinction is the central analytic step.

\begin{theorem}[Auxiliary-parameter inverse]\label{thm:parameter}
Let $F_0$ be analytic near a point $x_*$ with $F_0'(x_*)\ne0$ and let
$\phi=F_0^{-1}$ be the corresponding local inverse near $y_*=F_0(x_*)$.
Let $h$ be analytic near $x_*$. The branch of
\[
 F_0(x)+\varepsilon h(x)=y
\]
that equals $\phi(y)$ at $\varepsilon=0$ is analytic in $\varepsilon$ near
zero and satisfies
\begin{equation}\label{eq:param-formula}
 x=\phi(y)+\sum_{N\ge1}\frac{(-\varepsilon)^N}{N!}
 \left(\frac{\dd}{\dd y}\right)^{N-1}
 \left[\phi'(y)h(\phi(y))^N\right].
\end{equation}
\end{theorem}
\begin{proof}
Set $v=F_0(x)$ and $H(v)=h(\phi(v))$. The equation becomes
$v+\varepsilon H(v)=y$, and the required answer is $\phi(v)$. Choose a
small positively oriented circle $\Gamma$ around $y$ on which all
functions are analytic and $|\varepsilon H(z)|<|z-y|$. Rouch\'e's theorem
gives one zero $v$ inside the circle. By the residue theorem,
\[
 \phi(v)=\frac{1}{2\pi i}\int_\Gamma
  \phi(z)\frac{1+\varepsilon H'(z)}{z-y+\varepsilon H(z)}\,\dd z.
\]
Expand the denominator geometrically, uniformly on the circle. For
$N\ge1$, the coefficient of $\varepsilon^N$ is
\[
 \frac{(-1)^N}{2\pi i}\int_\Gamma
 \left[\frac{\phi H^N}{(z-y)^{N+1}}
       -\frac{\phi H^{N-1}H'}{(z-y)^N}\right]\dd z.
\]
The integral of the derivative of $\phi H^N/(z-y)^N$ vanishes, so this is
\[
 \frac{(-1)^N}{N}\frac{1}{2\pi i}\int_\Gamma
   \frac{\phi'(z)H(z)^N}{(z-y)^N}\,\dd z
 =\frac{(-1)^N}{N!}\partial_y^{N-1}(\phi'H^N)(y).
\]
The constant coefficient is $\phi(y)$. This proves the formula and its
local analytic convergence. Analytic dependence also follows from the
implicit function theorem.
\end{proof}

For $F_0(x)=x$, the formula reduces to
\begin{equation}\label{eq:near-identity}
 g(y)=y+\sum_{N\ge1}\frac{(-1)^N}{N!}
 \frac{\dd^{N-1}}{\dd y^{N-1}}h(y)^N,
\end{equation}
provided evaluation at $\varepsilon=1$ is justified. Neither
\cref{thm:parameter} alone nor the existence of all derivatives at positive
$y$ is enough to justify that substitution uniformly as $y\to0$.
The next two sections supply the missing control for the finite class
\eqref{eq:model}.

\subsection{Relation to the supplied repository}
The companion volume \emph{Combinatorial Transseries and Their Inverses}
in the supplied repository states the phase-coordinate version of
\eqref{eq:param-formula}, a scaled Lagrange construction, and
residual-based inverse error transport \cite{companion}. Its discussion
of action collisions is the exponential counterpart of the power-weight
collisions here. The canonical volume's README explicitly distinguishes
operator and coefficient forms of Lagrange inversion and scopes its
formalization claims \cite{repo-readme}.

The present article specializes that architecture to a zero-endpoint
power--logarithmic input class. Its practical additions are the polynomial
operator formula below, the finite-border logarithmic remainder theorem,
and a concrete input/parser/truncation contract with two implemented
engines. The underlying Lagrange principle is classical; no priority claim
for the general inversion identity is intended. The repository's relevant
companion sections and README were consulted, not its entire large
formalization corpus. No Lean verification of this package is asserted.

\section{A closed coefficient formula for every block}
\label{sec:coefficients}

For a multi-index $\nn=(n_1,\ldots,n_r)$ define
\[
 N=|\nn|=\sum_j n_j,\qquad
 d=\nn\cdot\del,\qquad
 \nn!=\prod_j n_j!,\qquad
 Q_{\nn}(L)=\prod_jP_j(L)^{n_j}.
\]
An empty product is $1$.

\begin{theorem}[All-order power--log inverse]\label{thm:master}
For \eqref{eq:model}, define $C_{\boldsymbol0}(L)=1$. For $N\ge1$ set
\begin{equation}\label{eq:master}
 \boxed{\displaystyle
 C_{\nn}(L)=\frac{(-1)^N}{p^N\nn!}
       \prod_{k=1}^{N-1}\bigl(1+d+pk+\DL\bigr)Q_{\nn}(L).}
\end{equation}
Then the positive inverse has the generalized expansion
\begin{equation}\label{eq:master-expansion}
 \boxed{\displaystyle
 g(y)=s\sum_{\nn\in\N^r}s^{\nn\cdot\del}C_{\nn}(\log s),
 \qquad s=(y/c)^{1/p}.}
\end{equation}
For sufficiently small positive $s$ the sum converges absolutely.
At a fixed weight $d$, all $C_{\nn}$ with $\nn\cdot\del=d$ are added.
The degree bound is
\begin{equation}\label{eq:degree-bound}
 \deg C_{\nn}\le\sum_j n_j\deg P_j.
\end{equation}
\end{theorem}
\begin{proof}[Proof of the coefficient identity]
Temporarily replace the perturbation by
$\sum_j\varepsilon_j h_j(x)$ with
$h_j(x)=c x^{p+\delta_j}P_j(\log x)$. Apply
\cref{thm:parameter} with $F_0(x)=c x^p$ and expand the $N$th power by the
multinomial theorem. The coefficient of $\boldsymbol\varepsilon^{\nn}$ is
\begin{equation}\label{eq:before-operator}
 \frac{(-1)^N}{\nn!}\left(\frac{\dd}{\dd y}\right)^{N-1}
 \left[\phi'(y)\prod_j h_j(\phi(y))^{n_j}\right],
 \qquad \phi(y)=s.
\end{equation}
Now
\[
 \frac{\dd}{\dd y}=\frac{1}{cp}s^{1-p}\frac{\dd}{\dd s},
 \qquad
 \phi'(y)\prod_j h_j(\phi(y))^{n_j}
 =\frac{c^{N-1}}p s^{p(N-1)+d+1}Q_{\nn}(L).
\]
Each differentiation lowers the $s$-power by $p$ and acts on the
polynomial by that current power plus $\DL$. After $N-1$ differentiations,
the scalar factor is $p^{-N}$, the remaining power is $s^{d+1}$, and the
operator factors are $1+d+pk+\DL$ for $k=1,\ldots,N-1$. This proves
\eqref{eq:master}. Differentiation of a polynomial cannot increase its
degree, giving \eqref{eq:degree-bound}. The convergence and equality to
the selected inverse are established in \cref{thm:convergence} below.
\end{proof}

\subsection{Low orders and nonlogarithmic specialization}
For a single generator $x^\delta P(\log x)$, the first corrections are
\begin{align}
 g(y)={}&s-\frac1p s^{1+\delta}P(L)\notag\\
 &+\frac{s^{1+2\delta}}{2p^2}
       (1+2\delta+p+\DL)P(L)^2\notag\\
 &-\frac{s^{1+3\delta}}{6p^3}
       (1+3\delta+p+\DL)(1+3\delta+2p+\DL)P(L)^3+\cdots.
 \label{eq:first-general}
\end{align}
For two distinct generators, the mixed second-order coefficient is
\begin{equation}\label{eq:mixed-second}
 C_{\ee_i+\ee_j}(L)=\frac1{p^2}
   (1+\delta_i+\delta_j+p+\DL)P_i(L)P_j(L).
\end{equation}
There is no factor $1/2$ here, because $(\ee_i+\ee_j)!=1$.

When all $P_j$ are constants $a_j$, the operators become scalars:
\begin{align}
 C_{\nn}
 &=\frac{(-1)^N}{p\nn!}
     \rise{\frac{1+d}{p}+1}{N-1}\prod_j a_j^{n_j}\notag\\
 &=\frac1{1+d}\binom{-(1+d)/p}{N}
      \frac{N!}{\nn!}\prod_j a_j^{n_j}.
 \label{eq:nonlog-master}
\end{align}
Here $(a)^{\overline{k}}=a(a+1)\cdots(a+k-1)$ denotes a rising
factorial, with $(a)^{\overline0}=1$.
This is the generalized-power analogue of the classical coefficient
formula. The derivative operators in \eqref{eq:master} are exactly the
extra structure needed when logarithms occur.

\section{Why the generalized series actually converges}
\label{sec:convergence}

Write $m_j=\deg P_j$ and $M_L=1+|L|$. The relevant small quantities are
\begin{equation}\label{eq:small-variables}
 u_j=s^{\delta_j}M_L^{m_j},\qquad
 \eta(s)=\sum_j u_j.
\end{equation}
Each tends to zero, even though $|L|\to\infty$.

\begin{theorem}[Uniformly normalized analytic inversion]\label{thm:convergence}
There are constants $B>0$, $\rho>0$, and $s_0>0$, depending on the fixed
model, such that
\begin{equation}\label{eq:cauchy-bound}
 |C_{\nn}(L)|\le B\rho^{-|\nn|}M_L^{\sum_jm_jn_j}
 \qquad(L\le\log s_0).
\end{equation}
The series \eqref{eq:master-expansion} converges absolutely for all
sufficiently small positive $s$ and equals the branch in
\cref{thm:branch}. A total-degree truncation satisfies, for fixed $K$,
\begin{equation}\label{eq:totaldegree-remainder}
 g(y)-s\sum_{|\nn|\le K}s^{\nn\cdot\del}C_{\nn}(L)
 =\OO\!\left(s\eta(s)^{K+1}\right).
\end{equation}
\end{theorem}
\begin{proof}
Keep $L$ fixed temporarily and write $x=s w$. Introduce independent
complex parameters $\boldsymbol z$ in
\begin{equation}\label{eq:normalized-equation}
 w^p\left(1+\sum_j z_j w^{\delta_j}P_j(L+\log w)\right)-1=0.
\end{equation}
Choose a disk $|w-1|<r_0<1$ small enough that $w^p-1$ has exactly one zero
there, namely the simple zero $w=1$, and no zero on its boundary. The powers
and logarithm use their analytic branches on this disk. Such a disk exists
by the nonzero derivative $p$ at $1$.

Normalize the parameters by $v_j=z_jM_L^{m_j}$. On the closed disk,
$\log w$ is bounded. Since $P_j$ is a fixed polynomial,
\[
 \sup_{L\le\log s_0,\ |w-1|=r_0}
 \frac{|w^{p+\delta_j}P_j(L+\log w)|}{M_L^{m_j}}<\infty.
\]
Also $b=\min_{|w-1|=r_0}|w^p-1|>0$. Choose a positive polydisk radius
$\rho$ so that the sum of these bounds times $\rho$ is smaller than $b$.
Rouch\'e's theorem now gives exactly one zero of
\eqref{eq:normalized-equation} inside the disk, for every parameter point
in that polydisk, uniformly in $L$. A zero counted with multiplicity one
is simple. Local holomorphic dependence patches to a single holomorphic
function because the root in the disk is unique.

Its absolute value is bounded by $1+r_0$. Cauchy's estimates in a slightly
smaller polydisk give a bound $B\rho^{-|\nn|}$ for its coefficient in
$\boldsymbol v^{\nn}$. Returning to $\boldsymbol z$ multiplies the bound by
$M_L^{\sum m_jn_j}$, proving \eqref{eq:cauchy-bound}. The coefficient is
precisely $C_{\nn}(L)$ from \cref{thm:master}, by the local parameter identity.

At the actual values $z_j=s^{\delta_j}$, the normalized variables $v_j$
are the small quantities \eqref{eq:small-variables}. For sufficiently
small $s$, each $v_j<\rho/2$, so the multivariate series converges
absolutely. Its real root satisfies $w\to1$ and $x=s w>0$, hence is the
unique inverse germ already selected. For the tail of total degree $K+1$,
use
\[
 \sum_{|\nn|=N}\prod_j u_j^{n_j}
 \le\left(\sum_j u_j\right)^N
\]
and sum the geometric majorant. This gives
$B s(\eta/\rho)^{K+1}/(1-\eta/\rho)$, which implies
\eqref{eq:totaldegree-remainder}.
\end{proof}

The theorem is stronger than a purely formal reversion statement, but its
scope must be kept clear. The forward model is a \emph{finite exact} sum
with fixed positive gaps. A divergent infinite asymptotic input, a flat
remainder, or a parameter limit in which a gap tends to zero requires
additional arguments. Convergence here is not an assertion about every
transseries or every inverse problem.

\section{Weighted truncation and a finite-border remainder}
\label{sec:remainder}

Total perturbation degree is often the wrong output order. For example,
four copies of a small gap can have smaller total weight than two copies
of a larger one. The implementation therefore retains $\IA$ from
\eqref{eq:ideal}. Its finite \emph{exit border} is
\begin{equation}\label{eq:border}
 \BA=\{\nn+\ee_j:\nn\in\IA,
                  (\nn+\ee_j)\cdot\del\ge A,\ 1\le j\le r\}.
\end{equation}
Repeated border indices are removed. This border need not be a minimal
antichain; minimality is unnecessary for the estimate.

\begin{theorem}[Logarithm-aware weighted remainder]\label{thm:border}
Suppose $r>0$ and $A>0$. Define
\begin{equation}\label{eq:border-data}
 d_* =\min_{\boldsymbol b\in\BA}\boldsymbol b\cdot\del,
 \qquad
 k_* =\max_{\substack{\boldsymbol b\in\BA\\
               \boldsymbol b\cdot\del=d_*}}
                  \sum_j b_jm_j.
\end{equation}
Then
\begin{equation}\label{eq:border-remainder}
 g(y)=s\sum_{\nn\in\IA}s^{\nn\cdot\del}C_{\nn}(L)
       +\OO\!\left(s^{1+d_*}(1+|L|)^{k_*}\right).
\end{equation}
In the output variable, the remainder power is $(1+d_*)/p$. The bound is
safe even if some omitted coefficients cancel; it need not be sharp after
such cancellations.
\end{theorem}
\begin{proof}
Every $\nn\notin\IA$ dominates at least one border index componentwise.
Indeed, take any path from $\boldsymbol0$ to $\nn$ adding one unit vector
at a time, and consider its first exit from $\IA$. This exit lies in
$\BA$ and is componentwise no larger than $\nn$.

Set $v_j=\rho^{-1}s^{\delta_j}M_L^{m_j}$, with $\rho$ from
\cref{thm:convergence}. Its coefficient bound and the preceding cover imply
\begin{align*}
 \left|g-s\sum_{\nn\in\IA}s^{\nn\cdot\del}C_{\nn}(L)\right|
 &\le Bs\sum_{\nn\notin\IA}\boldsymbol v^{\nn}\\
 &\le Bs\sum_{\boldsymbol b\in\BA}\boldsymbol v^{\boldsymbol b}
           \prod_{j=1}^r\frac1{1-v_j}.
\end{align*}
For small $s$ the product of denominators is bounded. The remaining border
sum is finite. Among its terms with smallest power $d_*$, the largest
logarithmic degree is $k_*$. Every term with larger power is bounded by a
constant times $s^{d_*}M_L^{k_*}$ by \cref{lem:dominance}. This proves the
result.
\end{proof}

\subsection{Strict cutoff semantics}
To retain inverse powers $y^\lambda$ with $\lambda<q$, set
\begin{equation}\label{eq:cutoff-map}
 A=pq-1.
\end{equation}
Then $(1+d)/p<q$ is equivalent to $d<A$. The cutoff must satisfy
$q>1/p$ so that the leading block is present. The first possible omitted
power can be strictly larger than $q$ if there is a gap in the support.
The package reports the border power rather than degrading it to $q$.

The theorem also explains why appending \code{O[y]\^{}q} to an ordinary
expression is inappropriate in general. A block of power $q$ multiplied
by an unbounded logarithmic polynomial is not $\OO(y^q)$. The output
association has separate fields for its power, logarithmic degree, and
scale. The factor $1+|\log s|$ makes the bound nonnegative without relying
on the sign of $\log s$.

\subsection{Complexity and exact ordering}
The index count is at most
\[
 \prod_{j=1}^r\left(1+\left\lfloor\frac A{\delta_j}\right\rfloor\right),
\]
although most indices in this box are not retained. Each coefficient is
formed from a finite polynomial and $N-1$ first-order differential
operators. The exit border has at most $r|\IA|$ entries before
 deduplication. Collisions can greatly reduce the number of final blocks.

The package enumerates using exact comparisons, not floating floors of
approximate radical quotients. Exact algebraic constants are the primary
intended exponents. Other exact numeric real constants may be accepted
when every required comparison is proved by the kernel; an unresolved
comparison produces a failure, never a guessed ordering. Symbolic
exponents are deliberately excluded, because the support order and number
of retained indices can change with the parameter.

\section{A second engine: Newton inversion in finite jets}
\label{sec:newton}

The coefficient formula is nonrecursive and particularly transparent for
small support sets. A second construction is useful both computationally
and as an independent algebraic check. It works entirely in the finite
quotient obtained by discarding weights at least $A$.

Put $g(y)=s(1+U)$, where $U$ has positive weight. The normalized inverse
 equation is
\begin{equation}\label{eq:newton-G}
 G(U)=(1+U)^p+
 \sum_j s^{\delta_j}(1+U)^{p+\delta_j}
 P_j\bigl(L+\log(1+U)\bigr)-1=0.
\end{equation}
Its derivative with respect to $U$ is
\begin{align}
 G'(U)={}&p(1+U)^{p-1}\notag\\
 &+\sum_j s^{\delta_j}(1+U)^{p+\delta_j-1}
 \bigl((p+\delta_j)P_j+P_j'\bigr)
       \bigl(L+\log(1+U)\bigr).
 \label{eq:newton-Gprime}
\end{align}
Its leading coefficient is the nonzero scalar $p$, so it is a unit.

\begin{theorem}[Finite Newton reversion]\label{thm:newton}
Starting with $U_0=0$, perform
\begin{equation}\label{eq:newton-step}
 U_{k+1}=U_k-\frac{G(U_k)}{G'(U_k)},
\end{equation}
truncating all intermediate jets at weight $A$. The iteration terminates
with the unique positive-order solution of \eqref{eq:newton-G} in this
quotient. It equals the weighted truncation of \cref{thm:master}.
If the current error has valuation at least $v>0$, the next error has
valuation at least $2v$.
\end{theorem}
\begin{proof}
Let $\delta_{\min}$ be the smallest positive generator. Existence of a
formal solution follows either from the coefficient formula, or
coefficient by coefficient: the unknown coefficient of weight $d$ enters
$G(U)$ as $p$ times that coefficient, while every other contribution
involves only smaller weights. The finite ordering of weights below $A$
makes this induction legitimate. It also proves uniqueness.

Let $U_*$ be that solution and $E=U-U_*$. Formal Taylor expansion, valid
because $E$ has positive valuation, gives
\[
 G(U)=G'(U)E+\OO_{\mathrm{formal}}(E^2).
\]
Here $\OO_{\mathrm{formal}}(E^2)$ denotes membership in the ideal of
valuation at least $2\val E$, not an analytic big-$O$ assertion in $s$.
Dividing by the unit $G'(U)$ preserves valuation and proves doubling.
Initially $\val U_*\ge\delta_{\min}$, unless the solution is identically
zero. Thus after finitely many steps the error has valuation at least
$A$ and vanishes in the quotient. The exact coefficient construction
satisfies the same equation, so uniqueness identifies the two answers.
If there are no generators, $U_*=0$ from the outset.
\end{proof}

For a sparse-jet implementation one need not manipulate expressions
containing unevaluated logarithms of series. The following recurrence is
more reliable.

\begin{lemma}[Composition coefficients]\label{lem:composition}
For a scalar $a$ and a polynomial $P$, define
\[
 Q_0(L)=P(L),\qquad
 Q_{k+1}(L)=\frac{(a-k)Q_k(L)+Q_k'(L)}{k+1}.
\]
Then, formally in $t$,
\begin{equation}\label{eq:composition}
 (1+t)^a P(L+\log(1+t))=\sum_{k\ge0}Q_k(L)t^k.
\end{equation}
\end{lemma}
\begin{proof}
Let $H(t,L)$ denote the left side. Direct differentiation gives
$(1+t)\partial_tH=aH+\partial_LH$. Comparing the coefficient of $t^k$
gives $(k+1)Q_{k+1}+kQ_k=aQ_k+Q_k'$, which is the recurrence.
The constant coefficient is $P(L)$, so all coefficients agree.
\end{proof}

The algorithm stores a jet as a sorted list of pairs
$\{d,Q_d(L)\}$. Addition merges equal weights; multiplication adds
weights and multiplies polynomials; reciprocal uses a finite geometric
series after dividing by the constant coefficient. For composition,
\cref{lem:composition} generates $Q_k$ while powers of $U$ are multiplied
until none remain below the current cutoff. Every operation terminates.

A mathematical operation count for Newton need not translate directly to
a speedup in this research implementation. The code uses simple sparse
lists and repeated exact simplification, prioritizing auditability over
large-order performance. No performance comparison between the two
Wolfram implementations was run. A production implementation could cache
polynomial products, use a canonical number-field representation for
algebraic weights, and increase the cutoff progressively during Newton
iteration instead of computing at full precision from the start.

\section{Residuals, actual inverse errors, and certification}
\label{sec:residual}

A correct coefficient calculation must cancel the forward equation, but
one must still relate that cancellation to the inverse function. The
correct scaling depends on $p$.

\begin{theorem}[Residual transport near a monomial core]
\label{thm:residual}
Let $f$ satisfy \eqref{eq:model}, let $g$ be its positive inverse, and let
$a(y)/s\to1$. If
\[
 f(a(y))-y=\OO\bigl(s^{p+v}(1+|\log s|)^k\bigr),
\]
then
\begin{equation}\label{eq:residual-transport}
 a(y)-g(y)=\OO\bigl(s^{1+v}(1+|\log s|)^k\bigr).
\end{equation}
More generally, an absolute forward residual $R(y)$ gives inverse error
$\OO(s^{1-p}|R(y)|)$ whenever the indicated estimate holds on the
segment joining $a(y)$ and $g(y)$.
\end{theorem}
\begin{proof}
Both endpoints are asymptotic to $s$, so that segment lies in
$[s/2,2s]$ for sufficiently small $y$. The derivative formula in the proof
of \cref{thm:branch} gives constants $b_1,b_2>0$ such that
\[
 b_1 s^{p-1}\le f'(x)\le b_2 s^{p-1}\qquad(s/2\le x\le2s).
\]
The mean-value theorem gives
$|a-g|\le |f(a)-f(g)|/(b_1s^{p-1})$, proving the claims.
\end{proof}

For a requested inverse cutoff $q$, the normalized cutoff is $A=pq-1$.
The corresponding forward residual cutoff in powers of $y$ is
\begin{equation}\label{eq:residual-cutoff}
 r=1+\frac A p=q+\frac{p-1}{p}.
\end{equation}
Thus ``residual zero below $q$'' is not the correct generic test when
$p\ne1$. The package computes the normalized residual
$f(s(1+U))/(cs^p)-1$ below weight $A$, and reports its corresponding
cutoff \eqref{eq:residual-cutoff}.

An asymptotic scale is not an explicit numerical error bar. For a
specified $y$, a certificate needs an actual neighborhood, derivative
bound, and residual bound. The following elementary result states exactly
what is required.

\begin{theorem}[Two-sided residual certificate]\label{thm:certificate}
Suppose $f$ is continuous on $[a-h,a+h]$, differentiable in its interior,
and $f'(x)\ge m>0$ there. If $|f(a)-y|\le mh$, there is a unique solution
$g$ of $f(g)=y$ in the interval, and
\[
 |g-a|\le\frac{|f(a)-y|}{m}.
\]
\end{theorem}
\begin{proof}
The mean-value theorem gives
$f(a-h)\le f(a)-mh\le y$ and
$f(a+h)\ge f(a)+mh\ge y$. Continuity gives a root between the endpoints,
and the positive derivative gives uniqueness. Applying the same derivative
bound between $a$ and that root proves the error estimate.
\end{proof}

The package does not search for an interval satisfying these hypotheses
or perform outward-rounded interval arithmetic. Accordingly, its field
\code{"NumericalIntervalCertificate"} is \code{False}. Its
\code{"FormalResidualZero"} field means exact cancellation in the finite
jet algebra of the supplied model. These are deliberately different
notions. The branch and asymptotic remainder theorems justify the symbolic
output; a numerical certificate at a particular argument remains an
additional step.

\section{Inverting a finite forward expansion with a remainder}
\label{sec:input-remainder}

One often knows a finite asymptotic model rather than an exact finite
function. This can be handled without pretending that the omitted input
terms are zero.

\begin{theorem}[Stability under a differentiated remainder]
\label{thm:input-remainder}
Let $f_0$ be a finite model \eqref{eq:model}, and suppose the actual
function is $F(x)=f_0(x)+r(x)$, with
\begin{equation}\label{eq:input-contract}
 r(x)=\OO\bigl(x^\beta(1+|\log x|)^k\bigr),\qquad
 r'(x)=\OO\bigl(x^{\beta-1}(1+|\log x|)^k\bigr),
 \qquad \beta>p,
\end{equation}
where $k$ is a nonnegative integer. Then $F$ has a positive inverse germ
$G$, and
\begin{equation}\label{eq:input-transport}
 G(y)-f_0^{-1}(y)
 =\OO\bigl(s^{\beta-p+1}(1+|\log s|)^k\bigr).
\end{equation}
Consequently, coefficients of inverse powers strictly below
\begin{equation}\label{eq:max-input-cutoff}
 q_{\max}=\frac{\beta-p+1}{p}
\end{equation}
are determined by the finite model. A remainder for a truncated model
inverse must be combined with \eqref{eq:input-transport}.
\end{theorem}
\begin{proof}
The value bound gives $F(x)/(cx^p)\to1$, and the derivative bound is
$o(x^{p-1})$. Thus $F'(x)\sim cp x^{p-1}>0$ and the same branch argument
as before gives $G(y)/s\to1$. At $a(y)=f_0^{-1}(y)\sim s$ the residual
$F(a)-y$ equals $r(a)$, which has size
$\OO(s^\beta(1+|\log s|)^k)$. Apply the derivative estimate for $F$ and
\cref{thm:residual}. This proves \eqref{eq:input-transport} and then
\eqref{eq:max-input-cutoff}.
\end{proof}

The theorem does not say that every coefficient at the boundary power is
undetermined in every particular problem. It gives a uniform safe rule:
all complete blocks strictly below the boundary are determined, while the
unknown remainder can affect the boundary block.

The derivative hypothesis cannot simply be omitted. The function
\[
 F(x)=x+x^2\sin(x^{-3})
\]
satisfies $F(x)=x+\OO(x^2)$ but
\[
 F'(x)=1+2x\sin(x^{-3})-3x^{-2}\cos(x^{-3})
\]
changes sign arbitrarily close to zero. The value approximation alone
therefore does not imply a single monotone inverse branch on a full
punctured neighborhood. Other hypotheses can replace the derivative
bound, for example a separate proved monotonicity and lower-slope
condition; they are not inferred automatically here.

The option
\begin{lstlisting}
"InputRemainder" -> {beta, k}
\end{lstlisting}
asserts both bounds in \eqref{eq:input-contract}. The package checks
$\beta>p$ and rejects a requested cutoff $q>q_{\max}$. It does \emph{not}
prove those input bounds. When combining errors, the smaller power wins;
at equal powers the larger logarithmic degree wins. For example,
\begin{lstlisting}
RealInverseSeries[x + x^2 (1 + Log[x]), {x, y}, 4,
  "InputRemainder" -> {4, 0}]
\end{lstlisting}
has a model truncation error $\OO(y^4(1+|\log y|)^3)$ and an input-induced
error $\OO(y^4)$. The former dominates, so the reported logarithmic degree
is $3$. Requesting cutoff $5$ from the same remainder information is
rejected rather than silently treating an unknown fourth-order forward
term as zero.

\section{The logarithmic example to arbitrary order}
\label{sec:log-example}

For $f(x)=x+x^2(1+\log x)$ we have $p=c=\delta=1$ and $P(L)=1+L$.
The all-order answer is
\begin{equation}\label{eq:log-all}
 g(y)=y+\sum_{n\ge1}y^{n+1}A_n(\log y),\qquad
 A_n(L)=\frac{(-1)^n}{n!}
       \prod_{k=1}^{n-1}(n+k+1+\DL)(1+L)^n.
\end{equation}
This is both a convergent generalized series for sufficiently small $y$
and a power--logarithmic asymptotic expansion to arbitrary order. Keeping
$n=1,\ldots,N$ gives
\begin{equation}\label{eq:log-tail}
 g(y)=y+\sum_{n=1}^N y^{n+1}A_n(\log y)
 +\OO\bigl(y^{N+2}(1+|\log y|)^{N+1}\bigr).
\end{equation}

The first five coefficient polynomials are
\begin{align}
 A_1(L)&=-L-1,\notag\\
 A_2(L)&=2L^2+5L+3,\notag\\
 A_3(L)&=-5L^3-\frac{41}{2}L^2-27L-\frac{23}{2},\notag\\
 A_4(L)&=14L^4+\frac{241}{3}L^3+\frac{335}{2}L^2
                       +151L+\frac{299}{6},\notag\\
 A_5(L)&=-42L^5-\frac{3727}{12}L^4-\frac{5365}{6}L^3
             -1256L^2-\frac{5177}{6}L-\frac{2789}{12}.
 \label{eq:log-polys}
\end{align}

\begin{proposition}[Catalan diagonal and a common factor]
\label{prop:catalan}
The coefficient of $L^n$ in $A_n$ is
\[
 (-1)^n\operatorname{Cat}_n
 =\frac{(-1)^n}{n+1}\binom{2n}{n}.
\]
Every $A_n$ is divisible by $1+L$.
\end{proposition}
\begin{proof}
Differentiation lowers polynomial degree, so only the scalar parts of
all operators contribute to $L^n$. Their product is
$(n+2)(n+3)\cdots(2n)=(2n)!/(n+1)!$. Multiplication by $(-1)^n/n!$
gives the first formula. There are at most $n-1$ derivatives applied to
$(1+L)^n$. Each derivative leaves at least one factor $1+L$, proving
 the second assertion.
\end{proof}

The common factor is also consistent with the exact fixed point:
$f(e^{-1})=e^{-1}$. This is a structural consistency check, not a proof
that the displayed series converges at $y=e^{-1}$; the local convergence
theorem only guarantees some sufficiently small interval.

The Catalan diagonal makes the logarithmic factor in \eqref{eq:log-tail}
sharp for each fixed $N$. Indeed, the first omitted block has nonzero
leading coefficient $(-1)^{N+1}\operatorname{Cat}_{N+1}$, and every later
block is smaller. Thus
\[
 \frac{|g(y)-g_N(y)|}
 {y^{N+2}|\log y|^{N+1}}\longrightarrow\operatorname{Cat}_{N+1}
 \quad(y\to0^+).
\]
This limit is a useful numerical diagnostic that tests more than just
small absolute error.

\section{The irrational-power example and its convergence radius}
\label{sec:irrational-example}

For the entire family $f(x)=x+x^\alpha$, $\alpha>1$, put
$\delta=\alpha-1$. The answer becomes
\begin{equation}\label{eq:alpha-all}
 g(y)=\sum_{n\ge0}b_n y^{1+n\delta},\qquad b_0=1,
 \qquad
 b_n=\frac{(-1)^n}{n!}\fall{\alpha n}{n-1}\quad(n\ge1).
\end{equation}
Equivalently,
\begin{equation}\label{eq:alpha-gamma}
 b_n=(-1)^n
 \frac{\Gamma(\alpha n+1)}{\Gamma(n+1)\Gamma(\delta n+2)}.
\end{equation}
The notation $(a)^{\underline{k}}=a(a-1)\cdots(a-k+1)$ denotes a
falling factorial, with the empty product equal to $1$.
Here the gamma quotient is merely a convenient rewrite of a finite
product; computing coefficients needs neither gamma evaluation nor a
special-function inverse.

\begin{proposition}[Radius in the natural small variable]
\label{prop:alpha-radius}
The power series $g(y)/y=\sum_{n\ge0}b_n z^n$ in
$z=y^\delta$ has radius
\begin{equation}\label{eq:alpha-radius}
 R_\alpha=\frac{(\alpha-1)^{\alpha-1}}{\alpha^\alpha}.
\end{equation}
In particular, the generalized series converges for positive $y$ with
$y^{\alpha-1}<R_\alpha$ and there equals the positive inverse.
\end{proposition}
\begin{proof}
The coefficient formula is \eqref{eq:master} with $p=1$, $P=1$ and
$d=n\delta$. Its scalar product is
$(\alpha n)(\alpha n-1)\cdots(\alpha n-n+2)$, giving
\eqref{eq:alpha-all}. Applying Stirling's formula \cite{stirling} to the three gamma
factors in \eqref{eq:alpha-gamma} gives
\[
 |b_n|\sim\frac{\sqrt\alpha}{\sqrt{2\pi}\,\delta^{3/2}}
 n^{-3/2}\left(\frac{\alpha^\alpha}{\delta^\delta}\right)^n.
\]
The root test proves the radius. On a small positive interval the sum is
already the inverse by \cref{thm:convergence}. To continue along
$0<z<R_\alpha$, its germ solves $w+zw^\alpha=1$ with $w=g(y)/y$.
This equation has exactly one positive solution and nonzero derivative
$1+\alpha z w^{\alpha-1}$. Along any compact positive subinterval that
solution stays bounded away from zero and continues analytically by the
implicit function theorem. Uniqueness of analytic continuation identifies
it with the convergent series throughout the indicated interval.
\end{proof}

For $\alpha=\sqrt2$, the first terms are
\begin{align}
 g(y)={}&y-y^{\sqrt2}+\sqrt2\,y^{2\sqrt2-1}
 -\frac{6-\sqrt2}{2}y^{3\sqrt2-2}\notag\\
 &+\frac{(4\sqrt2)(4\sqrt2-1)(4\sqrt2-2)}{24}
      y^{4\sqrt2-3}+\cdots.
 \label{eq:sqrt2-five}
\end{align}
Truncating immediately before the fourth correction yields exactly the
order requested in the question. The exponents are not rounded to a
rational grid. Since the coefficients are constants, the remainder has
logarithmic degree zero.

The elementary check $\alpha=2$ recovers
$g(y)=(\sqrt{1+4y}-1)/2$ and the alternating Catalan coefficients.
Its ordinary Taylor radius is $1/4$, consistent with
\eqref{eq:alpha-radius}. This is a useful limit case for both engines.

\section{Mixed generators, collisions, and a nonunit core}
\label{sec:mixed}

\subsection{A genuinely mixed scale}
Consider
\[
 f(x)=x+x^{\sqrt2}+x^2(1+\log x),\qquad
 \delta=\sqrt2-1.
\]
At inverse cutoff $q=3$, the normalized cutoff is $A=2$. The retained
multi-indices $(n,m)$ satisfy $n\delta+m<2$. There are eight of them,
listed here in increasing weight:
\begin{center}
\begin{tabular}{cccccccc}
\toprule
$(0,0)$&$(1,0)$&$(2,0)$&$(0,1)$&$(3,0)$&$(1,1)$&$(4,0)$&$(2,1)$\\
\midrule
$0$&$\delta$&$2\delta$&$1$&$3\delta$&$\delta+1$&$4\delta$&$2\delta+1$\\
\bottomrule
\end{tabular}
\end{center}
The full answer below that cutoff is
\begin{align}
 g(y)={}&y-y^{\sqrt2}+\sqrt2\,y^{2\sqrt2-1}
      -(1+L)y^2-\frac{6-\sqrt2}{2}y^{3\sqrt2-2}\notag\\
 &+\bigl((\sqrt2+2)(1+L)+1\bigr)y^{\sqrt2+1}
      +b_4y^{4\sqrt2-3}\notag\\
 &-\frac12\left[(2\sqrt2+1)(2\sqrt2+2)(1+L)
                         +4\sqrt2+3\right]y^{2\sqrt2}\notag\\
 &+\OO\bigl(y^3(1+|L|)^2\bigr),\qquad L=\log y,
 \label{eq:mixed-full}
\end{align}
where $b_4$ is the positive coefficient in \eqref{eq:sqrt2-five}.
The final displayed polynomial follows by applying two operators to
$1+L$ with $n=2,m=1$. The border weight is $2$, from $(0,2)$, and its
logarithmic degree is $2$. In contrast, a total-degree cutoff would
incorrectly omit some of the pure irrational-power terms that are larger
than retained mixed terms.

\subsection{Coincident weights must be added}
For
\[
 f(x)=x+x^{3/2}+x^2,
\]
the generators are $1/2$ and $1$. The weight $1$ receives two
contributions: $3/2$ from the multi-index $(2,0)$ and $-1$ from $(0,1)$.
Therefore
\[
 g(y)=y-y^{3/2}+\frac12y^2+\OO(y^{5/2}).
\]
Keeping the two contributions as separate ``ordered sectors'' would not
be a canonical expansion. The package groups them before presenting
blocks, and both algorithms pass this collision check in the independent
verification.

\subsection{The core coefficient and exponent both matter}
For
\[
 f(x)=3x^2\bigl(1+x^{1/2}(1+\log x)+2x\bigr)
\]
put $s=(y/3)^{1/2}$ and $L=\log s$. Then
\begin{align}
 g(y)={}&s-\frac12s^{3/2}(1+L)\notag\\
 &+s^2\left[\frac12(1+L)^2+\frac14(1+L)-1\right]
 +\OO\bigl(s^{5/2}(1+|L|)^3\bigr).
 \label{eq:nonunit-example}
\end{align}
The first contribution at $s^2$ comes from two copies of the half-power
generator, while $-s^2$ comes from the $2x$ generator. The powers of $p$
in the master formula are essential. Treating this as a near-identity
problem in $y$ without changing the derivative operator would give
incorrect coefficients and an incorrect residual cutoff.

\section{The Wolfram Language package}
\label{sec:package}

\subsection{Loading and the principal function}
Unpack the archive and load the package by its actual filesystem path:
\begin{lstlisting}
Get["/path/to/real-inverse-series/Kernel/RealInverseSeries.wl"];
Clear[x, y];
r = RealInverseSeries[x + x^2 (1 + Log[x]), {x, y}, 5];
r["Expansion"]
r["Remainder"]
\end{lstlisting}
The third argument is a \emph{strict power cutoff}, not a number of terms
and not a perturbation order. A request with cutoff $5$ includes the
complete blocks of powers $1,2,3,4$. For the irrational example the call
reproducing the question's four displayed terms is
\begin{lstlisting}
RealInverseSeries[x + x^Sqrt[2], {x, y}, 4 Sqrt[2] - 3]
\end{lstlisting}
All input and output variable symbols must be unassigned and distinct.
The forward expression must be independent of the output symbol: it
cannot simultaneously be used as an input parameter.

\subsection{Accepted input and deliberate failures}
The parser expands a finite sum of terms $x^aP(\log x)$, collects equal
powers, and proves the conditions on the smallest remaining power. Its
leading coefficient must be a positive scalar, not a polynomial in the
logarithm. Every power and the cutoff must be exact, finite, numeric real
constants whose required comparisons can be established exactly. Rational
numbers and exact algebraic numbers are the primary use cases. Constants
such as $\pi$ may also be accepted when all comparisons are decidable by
exact simplification. They are not converted to floating-point numbers.

Polynomial coefficients may contain symbolic parameters if their reality
is proved using \code{Assumptions}. The leading coefficient must likewise
be proved positive. For example,
\begin{lstlisting}
RealInverseSeries[x + a x^(3/2), {x, y}, 5/2,
  Assumptions -> Element[a, Reals]]
\end{lstlisting}
returns the polynomial $y-a y^{3/2}+\tfrac32a^2y^2$ with a conditional
remainder. No uniformity in unbounded parameter sets is asserted.

Inputs such as \code{x + x\^{}1.5}, a logarithmic leading core, a symbolic
exponent, or \code{x + Exp[-1/x]} return \code{Failure} rather than an
unsupported expansion. Inexact coefficients are also rejected: exact
symbolic residual cancellation is part of the contract. A user may apply
\code{N} to the returned expansion for numerical evaluation after the
symbolic construction.

\subsection{Options}
\begingroup\small
\begin{longtable}{@{}>{\raggedright\arraybackslash}p{.37\textwidth}>{\raggedright\arraybackslash}p{.58\textwidth}@{}}
\toprule
Option and default & Meaning\\\midrule\endhead
\code{Method -> "Lagrange"} & Use the closed multi-index formula.
\code{"Newton"} selects the independent sparse-jet Newton engine.\\[4pt]
\code{Assumptions -> True} & Conditions used to prove coefficient reality,
core positivity, equality, and exact comparison. They are assumptions,
not numerical evidence.\\[4pt]
\code{"InputRemainder" -> None} & Treat the supplied finite model as exact.
A pair \code{\{beta,k\}} instead asserts the two differentiated bounds
in \eqref{eq:input-contract}.\\[4pt]
\code{"MaxIndices" -> 10000} & Limit the enumerated order ideal. Exceeding
the limit returns a failure; no partial expansion is presented as complete.\\[4pt]
\code{"VerifyResidual" -> True} & Recompose the output in finite jets and
require exact cancellation below the appropriate residual cutoff.
\code{False} skips that check and marks the field as not computed.\\
\bottomrule
\end{longtable}
\endgroup

\subsection{Result fields and representation}
The principal return value is an association, not a bare expression.
The most useful fields are as follows.
\begingroup\small
\begin{longtable}{@{}>{\raggedright\arraybackslash}p{.40\textwidth}>{\raggedright\arraybackslash}p{.55\textwidth}@{}}
\toprule
Field & Meaning\\\midrule\endhead
\code{"Expansion"} & The finite expression, with positive-real powers of
$y/c$ and logarithms substituted.\\[4pt]
\code{"Remainder"} & An association containing the nonnegative error scale,
power, logarithmic degree, provenance, and the explicit statement that it
is not a numerical interval certificate. Scale $0$ denotes an exact
inverse of the supplied exact model.\\[4pt]
\code{"Blocks"} & Sorted pairs $\{\lambda,Q(\ell)\}$ representing
$(y/c)^\lambda Q(\log(y/c)/p)$. Equal powers are merged.\\[4pt]
\code{"LogVariable"}, \code{"LogSubstitution"} & The private polynomial
variable $\ell$ and its intended substitution. Its generated name is not
part of the interface.\\[4pt]
\code{"Core"}, \code{"Generators"} & The parsed $c,p$ and the pairs
$\{\delta_j,P_j(\ell)\}$.\\[4pt]
\code{"Cutoff"}, \code{"IndexCount"} & The requested strict inverse cutoff
and the number of retained multi-indices before collisions.\\[4pt]
\code{"FormalResidualZero"} & Whether the optional exact recomposition
check completed successfully.\\[4pt]
\code{"VerifiedResidualCutoff"} & The forward power cutoff
$q+(p-1)/p$, not generally $q$.\\[4pt]
\code{"Branch"}, \code{"Conditions"} & The positive inverse-germ
specification and supplied conditions.\\[4pt]
\code{"ForwardRemainderContract"} & Whether the finite input is exact or
which bounds the user has assumed.\\
\bottomrule
\end{longtable}
\endgroup
The internal sparse representation is retained for
\code{InverseResidual}; consumers should not edit \code{"InternalData"}.
This release does not define overloaded arithmetic on its result
associations. To compose or combine several results, their remainder
contracts must also be propagated; manipulating only
\code{"Expansion"} discards that information.

\subsection{Residual inspection and generic perturbation inversion}
For a result \code{r},
\begin{lstlisting}
InverseResidual[r]
InverseResidual[r, 6]
\end{lstlisting}
returns the formal expansion of the finite model's residual. The first
call uses \eqref{eq:residual-cutoff}. The second inspects all residual
powers strictly below $6$, and can reveal the first nonzero omitted
coefficient. It does not manufacture additional inverse coefficients.
For the two-block logarithmic approximation,
\[
 f\bigl(y-y^2(1+\log y)\bigr)-y
 =-y^3(2\log^2y+5\log y+3)+\OO(y^4(1+|\log y|)^3).
\]
The sign is opposite to the next inverse correction, as expected.

The separate utility
\begin{lstlisting}
PerturbativeInverse[phi, h, {x, y, eps}, n]
\end{lstlisting}
implements \eqref{eq:param-formula} through degree $n$ in
\code{eps}. Here \code{phi} is an expression in $y$ for a specified local
inverse of a core $F_0$, and \code{h} is the perturbation expressed in $x$.
For example,
\begin{lstlisting}
PerturbativeInverse[y, x^2 (1 + Log[x]), {x, y, eps}, 2]
\end{lstlisting}
gives
\[
 y-\varepsilon y^2(1+\log y)
 +\varepsilon^2y^3(2\log^2y+5\log y+3).
\]
This function is an exact finite differential formula, not an automatic
real-branch detector. A caller supplying an arbitrary \code{phi} must
ensure that it is the intended local inverse and satisfies the theorem's
hypotheses. Setting \code{eps -> 1} requires a convergence or asymptotic
argument; the principal power--log function supplies that argument only
for its supported model class.

\subsection{Relationship to built-in series operations}
The historical question demonstrates a particular attempt involving
\code{InverseFunction} and \code{Series}; it is not a benchmark of present
kernel behavior. Official documentation describes \code{InverseSeries}
as reversion of a series and describes the rational power-grid structure
of \code{SeriesData}. It also explicitly allows logarithms or slower
functions as coefficients in some series, and leaves symbolic powers
outside a \code{SeriesData} object \cite{wl-seriesdata,wl-inverseseries}.
Therefore it would be incorrect to say that built-in series have no
logarithmic support at all.

The reason for the custom representation is narrower: this algorithm
needs an arbitrary finitely generated real exponent support, complete
polynomial-logarithmic blocks, exact collision handling, and an explicit
analytic meaning for the remainder. It avoids reliance on how a
particular kernel version chooses to represent such a result. No
\code{PowerExpand} is used, so positive-real branch information is not
silently replaced by generally invalid complex identities.

\section{Verification record and reproducibility}
\label{sec:verification}

\subsection{What was executed}
The distribution contains an independent implementation in
\code{verification/reference.py}. It uses exact SymPy expressions for
weights and coefficients, and mpmath for numerical evaluation. Its two
engines implement the closed operator formula and Newton jet inversion
separately. The program \code{verification/verify.py} was executed and
recorded \textbf{79 successful checks} in
\path{verification/verification-report.json}. Its environment was Python
3.13.5, SymPy 1.14.0, and mpmath 1.3.0. These are recorded execution
versions, not minimum supported versions or claims about the latest
releases.

The symbolic checks include the direct derivative formula through eight
orders in the logarithmic example, the Catalan diagonal, explicit low
coefficients, irrational-power coefficients through eight orders, exact
cutoff boundaries, mixed supports, and resonant collision sums. For eight
models, the independent Newton and Lagrange engines agree and the forward
residual cancels at the required normalized order. The models include a
nonunit core, a fractional core exponent, logarithmic coefficients,
negative perturbations, and the exact monomial case. Five additional
nonunit-core coefficients were checked directly by differentiating in
the original core coordinate.

Numerical checks were performed with 250 decimal digits of working
precision. Roots were found by Newton iteration in the normalized
variable $w=x/s$, not by subtracting unscaled quantities near zero.
The normalized forward residual at each computed root was below
$10^{-240}$ in the recorded checks. A residual reported as zero in the
JSON file means zero at that finite working precision, not an exact
symbolic root. The numerical calculations use ordinary multiprecision
arithmetic, not interval arithmetic.

For the logarithmic example, \cref{tab:numeric} gives a small selection
of the recorded values. The scaled error is
\[
 \frac{|g(y)-g_N(y)|}{y^{N+2}|\log y|^{N+1}},
 \qquad g_N=y+\sum_{n=1}^N y^{n+1}A_n(\log y).
\]
The predicted limits are $2$ for $N=1$ and $14$ for $N=3$.
\begin{table}[htbp]
\centering
\begin{tabular}{@{}rrrl@{}}
\toprule
$y$ & corrections $N$ & $|g-g_N|$ & scaled error\\
\midrule
$10^{-2}$ & 1 & $2.420207\times10^{-5}$ & 1.141198\\
$10^{-2}$ & 3 & $1.486663\times10^{-7}$ & 3.305444\\
$10^{-4}$ & 1 & $1.268500\times10^{-10}$ & 1.495337\\
$10^{-4}$ & 3 & $5.096420\times10^{-16}$ & 7.082107\\
$10^{-8}$ & 1 & $5.895398\times10^{-22}$ & 1.737408\\
$10^{-8}$ & 3 & $1.163926\times10^{-34}$ & 10.108872\\
\bottomrule
\end{tabular}

\caption{Measured inverse errors at 250-digit working precision.
The table is generated by the verification program from its recorded
numerical values; entries are rounded for display.}
\label{tab:numeric}
\end{table}

To reproduce the independent checks from the distribution root, run
\begin{lstlisting}
python -m pip install -r verification/requirements.txt
python verification/verify.py
\end{lstlisting}
The second command regenerates the JSON report and the table source.
It requires no Wolfram kernel. The word ``independent'' here means a
separate implementation and a second algorithmic route, not an externally
reviewed or formally verified program.

\subsection{What was not executed}
The native file \code{Tests/RealInverseSeries.wlt} contains
\textbf{33 Wolfram Language regression tests}. They cover the two
question examples, logarithmic remainders, exact irrational cutoffs,
mixed weights, collisions, pure and nonunit cores, the two methods,
residual inspection, parameter assumptions, forward remainder limits,
resource limits, unsupported inputs, and the perturbation utility.
These tests were supplied but \emph{not run in a Wolfram kernel} during
preparation. Both attempts to access the connected Wolfram service
returned an unavailable-endpoint error, and no local kernel was present.
The independent Python results therefore must not be read as a native
Wolfram test pass.

The native command is
\begin{lstlisting}
wolframscript -file Tests/run-tests.wl
\end{lstlisting}
or, from a notebook after loading the package,
\begin{lstlisting}
TestReport["/path/to/real-inverse-series/Tests/RealInverseSeries.wlt"]
\end{lstlisting}
The runner records the actual kernel version and report when executed,
and exits unsuccessfully if the report does not succeed. It uses the
current documented \code{"ReportSucceeded"} property when available,
with an older-report count fallback \cite{wl-tests}. No native report
file is shipped as though such an execution had occurred.

The Wolfram source was manually reviewed and subjected to a delimiter,
string, and nested-comment balance check. That check can catch malformed
source structure, but it is not a Wolfram parser, evaluator, or proof of
semantic correctness. Native execution remains the next validation step
for adopting the package in a production workflow. The mathematical
proofs do not depend on that execution, but implementation defects are
not ruled out merely by proving the algorithm.

\subsection{Rebuilding the article}
From the \code{article} directory, run
\begin{lstlisting}
pdflatex -interaction=nonstopmode -halt-on-error real-inverse-series.tex
pdflatex -interaction=nonstopmode -halt-on-error real-inverse-series.tex
pdflatex -interaction=nonstopmode -halt-on-error real-inverse-series.tex
\end{lstlisting}
A small \code{build.sh} script in the distribution root performs these
passes. The TeX source uses standard AMS, NewTX, listings, and hyperref
packages. It includes the full Wolfram source from the adjacent
\code{Kernel} directory and the generated numerical table; these paths
are relative, so preserve the distribution layout when rebuilding.

\section{Boundaries and extensions of the theory}
\label{sec:extensions}

\subsection{Finite jets extracted from a more general forward function}
A forward expression need not itself be a finite power--log sum for its
inverse to be accessible to the theory. For example, an analytic factor
can be expanded to a finite Taylor polynomial with a differentiated
remainder. Combining that polynomial with finitely many power--log
terms produces a model to which \cref{thm:input-remainder} applies.
The truncation order must be selected from the inverse error law, not
merely copied from the desired inverse cutoff.

As a concrete example, let
\[
 F(x)=\sin x+x^2\log x
 =x+x^2\log x-\frac{x^3}{6}+\OO(x^5).
\]
Taylor's theorem also gives the derivative remainder $\OO(x^4)$.
Thus the finite model can be submitted as
\begin{lstlisting}
RealInverseSeries[x + x^2 Log[x] - x^3/6, {x, y}, 5,
  "InputRemainder" -> {5, 0}]
\end{lstlisting}
The main parser does not expand \code{Sin[x]} itself. This explicit
separation keeps the origin and order of the forward remainder visible.

\subsection{Other centers and the point at infinity}
For a right-hand inverse near $(x_0,y_0)$, first write
$t=x-x_0$ and $u=y-y_0$. If the translated positive-side equation is a
finite model in $t$, the package applies to that model, and the result
is translated back. The local side and the sign of the leading core
must be chosen explicitly; a left-hand branch may need a substitution
$t=x_0-x$ and a corresponding target sign change.

An infinity problem can sometimes be converted to an endpoint problem
by $t=1/x$ and a suitable target reciprocal. This is a mathematical
reduction, not an automatic feature of this release. The inverse
reconstruction may amplify errors, so a bound on $t$ must be transported
through $x=1/t$. In particular, a relative error in $t$ is preferable to
an unsupported absolute big-$O$ statement after reciprocation.

\subsection{Logarithmic leading cores and deeper transseries}
The input
\[
 f(x)=-x\log x
\]
is not in \eqref{eq:model}: the leading coefficient is not a constant.
Its small positive inverse is
\[
 g(y)=-\frac{y}{W_{-1}(-y)}.
\]
Indeed, writing $x=e^u$ gives $u e^u=-y$, and the branch $u\to-\infty$
selects $W_{-1}$ rather than $W_0$. Thus a logarithmic core changes the
natural scale and branch analysis; it cannot be accommodated by simply
allowing $p=1$ and a coefficient called ``$-\log x$'' in the present
unit algebra. The generic perturbation utility can operate around a
specified Lambert core at a fixed positive point, but it does not supply
an automatic all-order endpoint remainder theorem for that larger class.

Negative powers of $\log x$, nested logarithms, exponential sectors, and
oscillatory factors likewise require a larger coefficient algebra and
new composition estimates. General transseries provide a framework for
such extensions; Edgar's introduction gives broader context
\cite{edgar}. The present construction uses only finite generator
semigroups and polynomial logarithmic blocks, so its proofs and finite
algorithms do not depend on implementing the full transseries field.

A flat correction is particularly instructive. Replacing a finite
model by $f(x)+e^{-1/x}$ can leave every algebraic-logarithmic asymptotic
coefficient unchanged, yet change the actual inverse by a nonzero
beyond-all-orders term. The differentiated-remainder theorem controls
it at every fixed finite algebraic order, but an algebraic expansion
alone does not identify that exponentially small term. A package claiming
the exact finite model's convergent inverse for the perturbed function
would be making a false identification.

\subsection{Parameter limits and large-order computation}
All statements above are for fixed positive $p,c,\delta_j$ and fixed
polynomials. As $\delta_{\min}\to0^+$, the number of indices below a
fixed cutoff can diverge and $s^{\delta_{\min}}$ need not be small in a
joint limit. A uniform double-scaling theory would need a new small
parameter and uniform hypotheses. This is why the implemented parser
requires fixed exact exponents instead of attempting symbolic ordering
across parameter regions.

For large orders, the most promising engineering extensions are canonical
algebraic-number coordinates, memoized coefficient operators, sparse
polynomial arrays, a heap-based support enumerator, and precision-doubling
Newton truncation. A sharper optional remainder could inspect cancellation
in successive omitted blocks. The present border bound deliberately
avoids that unbounded search: it always terminates and remains valid even
when its first possible omitted block vanishes.

\section{Conclusion}

The two original examples are not exceptions requiring unrelated tricks.
They are instances of a single power--logarithmic inversion theorem.
The correct procedure is to select the positive real germ, normalize the
monomial core, use an auxiliary parameter at positive points, and only
then reorganize the coefficients by their actual asymptotic weights.
The logarithmic derivative operator produces every polynomial coefficient;
local finiteness makes irrational supports computable; normalized
analytic inversion supplies convergence for finite exact models; and a
finite exit border supplies a safe logarithmic remainder at arbitrary
power cutoffs.

The package implements this pipeline in two ways and separates the finite
expression, formal residual, input assumptions, and analytic error scale.
For the question's first example, it also corrects the ordinary-Landau
interpretation of the two-term remainder. For the second, it reproduces
the irrational exponents and coefficients exactly. The mathematical
results are proved here, independent symbolic and numerical checks are
recorded, and native Wolfram execution is explicitly left unclaimed.

\appendix
\section{Distribution map and source provenance}

\begingroup\small
\begin{longtable}{@{}>{\raggedright\arraybackslash}p{.49\textwidth}>{\raggedright\arraybackslash}p{.46\textwidth}@{}}
\toprule
Path & Contents\\\midrule\endhead
\code{article/real-inverse-series.tex} & This article's source.\\
\code{article/real-inverse-series.pdf} & Compiled article.\\
\code{article/numerical-table.tex} & Table generated by the numerical verification.\\
\code{Kernel/RealInverseSeries.wl} & Complete Wolfram package.\\
\code{Kernel/init.m} & Loader for the adjacent package file.\\
\code{Examples/examples.wl} & Reproducible usage examples.\\
\code{Tests/RealInverseSeries.wlt} & 33 native regression tests, not executed here.\\
\code{Tests/run-tests.wl} & Native runner that records the kernel version.\\
\code{verification/reference.py} & Independent exact reference implementation.\\
\code{verification/verify.py} & The executed verification program.\\
\path{verification/verification-report.json} & All 79 checks, coefficients,
Newton traces, numerical values, and execution versions.\\
\code{verification/verification-log.txt} & Console output of the executed run.\\
\code{verification/static-wl-check.py} & Structural source check, not a kernel parser.\\
\code{README.md}, \code{build.sh}, \code{LICENSE} & Usage, build instructions,
and reuse terms.\\
\bottomrule
\end{longtable}
\endgroup

The user's uploaded archive was used as the local question source.
The relevant repository companion sections were read in the source blob
\code{05c56ce5ee60788a67f365110b62810550cba817}; the group README had blob
\code{516aaa30e84cb0036a00f803cb6e8e3026febb8e}. These are content identifiers,
not claims that the entire repository was audited. External materials
are cited below and are not redistributed as though they were authored
for this package.

\section{Complete Wolfram Language package}
\label{app:source}

The following listing is included directly from the distributed file,
so the article and package source cannot silently diverge during a
normal rebuild. The complete executable file is preferable to copying
code from the PDF, whose line wrapping is for reading only.

\lstinputlisting[basicstyle=\ttfamily\scriptsize]{../Kernel/RealInverseSeries.wl}

\begin{thebibliography}{99}

\bibitem{question}
Vladimir Reshetnikov,
\emph{How to get an asymptotic of the real-valued branch of the inverse function?}
Mathematica Stack Exchange, question 236367, December 2020.
\href{https://mathematica.stackexchange.com/questions/236367/how-to-get-an-asymptotic-of-the-real-valued-branch-of-the-inverse-function}{Public question and discussion}.
The uploaded Markdown/PDF archive supplies the same two examples.

\bibitem{dlmf}
NIST Digital Library of Mathematical Functions,
\emph{\S1.10(vii), Lagrange inversion theorem}.
\url{https://dlmf.nist.gov/1.10#vii}.
Consulted September 7, 2026.

\bibitem{gessel}
Ira M. Gessel,
\emph{Lagrange Inversion}, 2016.
\url{https://arxiv.org/abs/1609.05988}.
In particular, Theorem 2.1.1 gives the classical coefficient identity.

\bibitem{surya}
Erlang Surya and Lutz Warnke,
\emph{Lagrange Inversion Formula by Induction}, 2023.
\url{https://arxiv.org/abs/2305.17576}.
An elementary formal-series treatment complementary to the local analytic
argument used here.

\bibitem{companion}
Vladimir Reshetnikov, ProveIt repository,
\emph{Combinatorial Transseries and Their Inverses},
chapter ``Coefficient calculus and the inversion engines.''
\href{https://github.com/VladimirReshetnikov/ProveIt/blob/main/Analysis/FabiusFunction/docs/semi-formalized-research-frontiers/drafts/series-and-transseries/Combinatorial_Transseries_Inverses/Combinatorial_Transseries_Inverses.tex}{Companion source in the supplied series-and-transseries directory}.
Relevant phase-coordinate, convergence, collision, and residual sections
consulted September 7, 2026.

\bibitem{repo-readme}
Vladimir Reshetnikov, ProveIt repository,
\emph{Series and transseries}, group README.
\href{https://github.com/VladimirReshetnikov/ProveIt/blob/main/Analysis/FabiusFunction/docs/semi-formalized-research-frontiers/drafts/series-and-transseries/README.md}{Group overview, provenance, and scoped formalization inventory}.
Consulted September 7, 2026.

\bibitem{stirling}
NIST Digital Library of Mathematical Functions,
\emph{\S5.11, Asymptotic expansions of the gamma function}.
\url{https://dlmf.nist.gov/5.11}.
The leading Stirling approximation is used in the radius calculation.

\bibitem{edgar}
G. A. Edgar,
\emph{Transseries for beginners}, Real Analysis Exchange
\textbf{35} (2009/2010), 253--310.
\url{https://arxiv.org/abs/0801.4877}.

\bibitem{wl-seriesdata}
Wolfram Research,
\emph{SeriesData}, Wolfram Language Documentation.
\url{https://reference.wolfram.com/language/ref/SeriesData.html}.
Consulted September 7, 2026.

\bibitem{wl-inverseseries}
Wolfram Research,
\emph{InverseSeries}, Wolfram Language Documentation.
\url{https://reference.wolfram.com/language/ref/InverseSeries.html}.
Consulted September 7, 2026.

\bibitem{wl-tests}
Wolfram Research,
\emph{TestReportObject}, Wolfram Language Documentation.
\url{https://reference.wolfram.com/language/ref/TestReportObject.html}.
Consulted September 7, 2026.

\end{thebibliography}
\end{document}
